\documentclass{article}

\usepackage{microtype}
\usepackage{graphicx}
\usepackage{subfigure}
\usepackage{booktabs}
\usepackage{hyperref}

\newcommand{\theHalgorithm}{\arabic{algorithm}}

% \usepackage{icml2025}        % blind review
\usepackage[accepted]{icml2025} % camera-ready

\usepackage{amsmath}
\usepackage{amssymb}
\usepackage{mathtools}
\usepackage{amsthm}

\usepackage[capitalize,noabbrev]{cleveref}

\theoremstyle{plain}
\newtheorem{theorem}{Theorem}[section]
\newtheorem{proposition}[theorem]{Proposition}
\newtheorem{lemma}[theorem]{Lemma}
\newtheorem{corollary}[theorem]{Corollary}
\theoremstyle{definition}
\newtheorem{definition}[theorem]{Definition}
\newtheorem{assumption}[theorem]{Assumption}
\theoremstyle{remark}
\newtheorem{remark}[theorem]{Remark}

% \usepackage[disable,textsize=tiny]{todonotes}
\usepackage[textsize=tiny]{todonotes}

% TikZ for framework diagram
\usepackage{tikz}
\usepackage{tikz-cd}
\usetikzlibrary{arrows.meta,positioning,matrix,fit,backgrounds}

% REFINED [old]: "Short Running Title" → [new]: actual short title
\icmltitlerunning{Where Models Miss: GUI Action Grounding}

\begin{document}

\twocolumn[
% REFINED [old]: "Paper Title" → [new]: actual paper title
\icmltitle{Where Models Miss: A Cross-Platform Analysis of GUI Action Grounding}

\icmlsetsymbol{equal}{*}

\begin{icmlauthorlist}
\icmlauthor{Aliaa Gheis}{equal,inst1}
\icmlauthor{Menna Abdelbaset}{equal,inst1}
\icmlauthor{Yara Hisham}{equal,inst1}
\icmlauthor{Nouran Hany}{equal,inst1}
\icmlauthor{Sarah Elzayat}{equal,inst1}
\end{icmlauthorlist}

\icmlaffiliation{inst1}{Computer Engineering Department, Cairo University Faculty of Engineering, Egypt, Cairo}

\icmlcorrespondingauthor{Aliaa Gheis}{aliaagheis@gmail.com}

\icmlkeywords{Vision Language Models, GUI Automation, Grounding}

\vskip 0.3in
]

\printAffiliationsAndNotice{\icmlEqualContribution}

\begin{abstract}
%
\end{abstract}

\section{Introduction}

\section{Related Work}

% ==========================================================
% ==========================================================
% Methods
% ==========================================================
% ==========================================================
\section{Method}

Click2Act is an open, extendable evaluation framework for GUI agents. it supports two evaluation modes ---point grounding and action prediction--- and is designed to be modular: new agents and benchmarks are registered independently and paired at runtime via YAML configuration files as illustrated in Figure~\ref{fig:framework}.

\begin{figure}[t]
    \centering
    % Click2Act framework diagram — vertical layout for two-column paper
% Total width ≈ 5.5 cm → fits comfortably in one column

\definecolor{c2aPrimary}{HTML}{2E86AB}
\definecolor{c2aSecondary}{HTML}{A23B72}
\definecolor{c2aTertiary}{HTML}{F18F01}
\definecolor{c2aNeutral}{HTML}{6C757D}
\definecolor{c2aSuccess}{HTML}{3BB273}
\definecolor{c2aText}{HTML}{212529}
\definecolor{c2aImplText}{HTML}{4A154B} % A deep, highly readable plum color
\tikzset{
  fwbox/.style={
    rectangle, rounded corners=2pt,
    inner sep=4pt, font=\scriptsize\sffamily,
    text=c2aText, align=center,
  },
  abstract/.style ={fwbox, fill=c2aPrimary!18,   draw=c2aPrimary,   line width=0.7pt, minimum width=1.8cm, minimum height=0.5cm},
  groupbox/.style ={
    fwbox, 
    fill=c2aSecondary!18, 
    draw=c2aSecondary, 
    line width=0.7pt, 
    minimum width=2cm, 
    text=c2aImplText, 
    font=\scriptsize\ttfamily\mdseries% Medium font weight in monospace
  },
  regbox/.style   ={fwbox, fill=c2aNeutral!15,   draw=c2aNeutral,   line width=0.7pt, minimum width=5.0cm, minimum height=0.5cm},
  pipeline/.style ={fwbox, fill=c2aTertiary!20,  draw=c2aTertiary,  line width=0.7pt, minimum width=1.8cm, minimum height=0.5cm},
  config/.style   ={fwbox, fill=c2aNeutral!12,   draw=c2aNeutral,   line width=0.5pt, dashed, minimum width=1.5cm, minimum height=0.45cm},
  layerlabel/.style={font=\tiny\bfseries\sffamily, text=c2aNeutral},
  arr/.style  ={-{Stealth[length=4pt]}, c2aNeutral,   line width=0.6pt},
  inherit/.style={arr, dashed, c2aPrimary},
  data/.style ={-{Stealth[length=4pt]}, c2aSuccess,   line width=0.6pt},
}

\centering
\begin{tikzpicture}[node distance=0.0cm]

% Layer 0 label
\node[layerlabel] (lbl0) at (0, 0) {Base Classes};

% Row 0: abstract bases
\node[abstract] (agentBase) at (-1.35, -0.45) {\texttt{GUIAgent}};
\node[abstract] (benchBase) at ( 1.35, -0.45) {\texttt{Benchmark}};

% Separator
\draw[c2aNeutral!40, dashed] (-2.8, -0.85) -- (2.8, -0.85);

% Layer 1 label
\node[layerlabel] at (0, -1.05) {Implementations};

% Row 1: two grouped implementation boxes
% REFINED [old]: default text color in groupbox → [new]: explicit c2aText text color + improved spacing for visual clarity
\node[groupbox] (agents) at (-1.35, -2.05) {%
  \vspace{10pt}%
  Holo2:4B\\[1pt]
  Aguvis:7B\\[1pt]
  UI-Agile:3B\\[1pt]
  Aguvis+UI-Zoomer%
};
\node[groupbox] (benchmarks) at ( 1.35, -2.05) {%
  AndroidControl\\[1pt]
  ScreenSpot-Pro%
};

% Separator
\draw[c2aNeutral!40, dashed] (-2.8, -3.15) -- (2.8, -3.15);

% Layer 2 label
\node[layerlabel] at (0, -3.35) {Registry};

% Row 2: registry bar
\node[regbox] (registry) at (0, -3.85)
  {\texttt{registry}};

% Separator
\draw[c2aNeutral!40, dashed] (-2.8, -4.30) -- (2.8, -4.30);

\node[layerlabel] at (0, -4.50)
  {Execution Scripts};

% YAML config nodes between registry and scripts
\node[config] (yamlA) at (-1.0, -4.95) {\texttt{agent.yaml}};
\node[config] (yamlB) at ( 1.0, -4.95) {\texttt{benchmark.yaml}};

% Row 3: pipeline scripts
\node[pipeline] (grounder) at (-1.35, -5.70) {\texttt{grounder.py}};
\node[pipeline] (scorer)   at ( 1.35, -5.70) {\texttt{action\_scorer.py}};

% Separator
\draw[c2aNeutral!40, dashed] (-2.8, -6.15) -- (2.8, -6.15);

% Layer 4 label + output
\node[layerlabel] at (0, -6.35) {Output};

\node[fwbox,
  fill=c2aSuccess!18, draw=c2aSuccess, line width=0.7pt,
  minimum width=3.2cm, minimum height=0.5cm]
  (csv) at (0, -6.85) {Results CSV};

% ── Edges ────────────────────────────────────────────────────────────────────

% base → implementations
\draw[inherit] (agentBase.south) -- (agents.north);
\draw[inherit] (benchBase.south) -- (benchmarks.north);

% implementations → registry
\draw[arr] (agents.south)  -- (agents.south  |- registry.north);
\draw[arr] (benchmarks.south) -- (benchmarks.south |- registry.north);

% registry → yaml → pipeline
\draw[arr] (registry.south -| yamlA.north) -- (yamlA.north);
\draw[arr] (registry.south -| yamlB.north) -- (yamlB.north);
% yaml -> pipeline
\draw[arr] (yamlA.south) -- (grounder.north);
\draw[arr] (yamlB.south) -- (scorer.north);

% pipeline → csv
\draw[data] (grounder.south) -- (grounder.south |- csv.north);
\draw[data] (scorer.south)   -- (scorer.south |- csv.north);

% ── Legend ───────────────────────────────────────────────────────────────────
\matrix [
  below=0.40cm of csv,
  column sep=5pt, row sep=2pt,
  matrix of nodes,
  nodes={font=\tiny\sffamily, text=c2aText, inner sep=1pt},
] {
  \tikz\node[abstract, minimum width=0.35cm, minimum height=0.20cm]{}; &
    Base class &[3pt]
  \tikz\node[groupbox, minimum width=0.35cm, minimum height=0.20cm]{}; &
    Impl. &[3pt]
  \tikz\node[regbox,   minimum width=0.35cm, minimum height=0.20cm]{}; &
    Registry \\
  \tikz\node[pipeline, minimum width=0.35cm, minimum height=0.20cm]{}; &
    Script &[3pt]
  \tikz\node[config,   minimum width=0.35cm, minimum height=0.20cm]{}; &
    YAML config &[3pt]
  \tikz\node[fwbox, fill=c2aSuccess!18, draw=c2aSuccess, line width=0.7pt,
             minimum width=0.35cm, minimum height=0.20cm]{}; &
    Output \\
};

\end{tikzpicture}

    \caption{%
        \textbf{Click2Act framework overview.}
        Abstract base classes (\texttt{GUIAgent}, \texttt{Benchmark}) are subclassed
        by concrete implementations and resolved at runtime via a registry.
        Each execution script receives an \texttt{agent.yaml} and a
        \texttt{benchmark.yaml}, runs inference over benchmark samples, and writes
        results to a CSV for offline analysis.
        Dashed arrows = inheritance; solid arrows = data flow.
    }
    \label{fig:framework}
\end{figure}

\subsection{Evaluated Models}

We evaluate four small pure-vision VLMs, all used with  no additional fine-tuning. Table~\ref{tab:models} summarizes each model.
\begin{table}[h]
    \centering
    \caption{Models evaluated in this study.}
    \label{tab:models}
    \begin{tabular}{lc}
    \toprule
    \textbf{Model} & \textbf{Params} \\
    \midrule
    UI-AGILE-3B      & 3B \\
    Holo2-4B         & 4B \\
    AGUVIS-7B        & 7B \\
    AGUVIS+UI-Zoomer & 7B \\
    \bottomrule
    \end{tabular}
    \end{table}
    
\noindent
All evaluated models fall in the 3B--7B parameter range, targeting practical deployment on consumer devices like laptops and workstations without cloud infrastructure. Despite this appeal, how much accuracy they sacrifice compared to flagships such as UI-Venus-1.5-30B-A3B ~\cite{uivenus15} and
Holo2-235B-A22B~\cite{holo2}, and where exactly they fail, remains an open question.

\textbf{AGUVIS-7B}~\cite{aguvis} is a fully open-source autonomous GUI agent built on Qwen2-VL that supports all platforms. It operates in two stages: (1) a planning stage where the model generates a structured \textit{inner monologue} ---consisting of a textual description of the current screen, internal reasoning toward the high-level goal, and low-level action instructions---, and (2) a grounding stage where the inner monologue is translated into executable actions. It was one of the first small fully open-source vision-based GUI agents to demonstrate competitive performance on ScreenSpot.

\textbf{Holo2-4B}~\cite{holo2} is a fully open-source autonomous GUI agent built on Qwen3-VL that supports all platforms. It is trained in two stages: (1) a large-scale supervised fine-tuning (SFT) stage on screen localization and element grounding, and (2) an online reinforcement learning stage using Group Relative Policy Optimization (GRPO) to refine its multi-step reasoning and navigation policies.

\textbf{UI-Agile-7B}

\textbf{AGUVIS+UI-Zoomer}

\subsection{Expirmented Benchmarks}
\subsection{Evaluation Metrics}



\section{Experiments}

\section{Conclusion}



\section*{Acknowledgements}

\section*{Impact Statement}

\bibliography{paper}
\bibliographystyle{icml2025}
\cite{aguvis}
\cite{li2024effects}
\appendix
\onecolumn
\section{Appendix}




\end{document}